\section{Local probes, periodic gates and random discovery}
\label{sec:acquisition}
The finite separation lemma removes the need to compare complete pair
images during acquisition. We now describe which local geometric work
remains and derive a discovery probability from fresh trajectories, rather
than assuming that an external producer finds every needed branch.

\subsection{A local validity test}
A recorded trajectory supplies encountered patches, not their unknown
periodic identities. Near a normal bridge, let $D(s,u)$ be the distance
between points on two local unit-speed boundary charts. At the normal
contacts,
\begin{equation}\label{eq:normal-hessian}
 \nabla D=0,\qquad D''=
 \begin{pmatrix}g^{-1}+\kappa_0&-g^{-1}\\
                -g^{-1}&g^{-1}+\kappa_1\end{pmatrix}.
\end{equation}
Its smallest eigenvalue is at least $\min(\kappa_0,\kappa_1)$.
Uniform third-derivative bounds therefore give a common neighborhood
where the Hessian stays positive. The inverse-function theorem, or a
Newton contraction on a smaller ball, isolates its unique zero from
sufficiently nearby encountered contacts. Enclosures of derivatives give
certified finite-accuracy root enclosures. At a true zero the two supporting
lines face one another. Their supporting half-planes show that this local
critical segment is the global shortest segment between the two convex
bodies; complete boundary maps are unnecessary for that assertion.

Clearance against other bodies is tested only near the candidate segment.
Outside fixed endpoint collars, a uniform clear bridge has distance at
least $c_0>0$ from the solid. Sample the distance-to-solid function on a
mesh of spacing less than $c_0/4$, with error less than $c_0/8$.
This distance function is 1-Lipschitz. Certified positive lower values,
minus the mesh radius, certify a smaller unobstructed tube. The endpoint
collars use the local convex charts and the separation bound $d_0$.
This needs only finitely many local distance queries. If certification
fails, return no record, not an assumed clear bridge.

\begin{lemma}[Local construction of genuine record descriptors]
\label{lem:probe}
With the certified probe of Section~\ref{sec:setting}, a bounded-time
three-impact trajectory can be either rejected or assigned to a verified
normal return branch, with its gap, descriptor \eqref{eq:signature},
and intrinsic endpoint offsets computed to prescribed errors. The rule
uses only neighborhoods of its contacts and candidate bridge. Every
trajectory in a common sufficiently small protected return rectangle of
a uniformly clear bridge is assigned to that bridge. The rule supplies
no complete body, translation class, period or completeness assertion.
\end{lemma}
\begin{proof}
Use a finite collection of local chart enclosures around the first,
second and third impacts. First and third must lie in the same connected
source patch, and the opposing chart must be distinct. Apply the positive
Hessian test above to the first two patches, retaining only facing
normal roots with gap certified below $R_+$. Check the middle tube and
endpoint collars. Shrink to fixed nested subpatches on which the
stationary two-flight branch is unique and its physical inequalities
hold. Assign a trajectory only if all three impacts lie in those patches
and its local reflection equations select this branch. Finite refinement
of the measured charts, arclength integrals and derivatives computes the
requested data. Refinement is stopped at a prescribed finite budget, and an indeterminate
case returns no record. Thus borderline invalid inputs cannot cause an
unbounded search. Uncertified cases are rejected.

For a bridge of length at most $R_-$ with the stated margins, the
normal root lies in the interior of the common chart and has the uniform
positive Hessian \eqref{eq:normal-hessian}. The difference $R_+-R_-$
protects the cutoff test. The tube, twist and adjacent-flight tests have
strict slack. The derivative and separation bounds give a uniform finite
refinement budget sufficient on this protected set. Continuous dependence of the local trajectories supplies
a common smaller protected rectangle on which all tests pass after
finite refinement. The result follows. The finiteness and certified
accuracy of the local measurements are precisely the probe-access
assumption, not a consequence of unmarked collision data alone.
\end{proof}

Applying Proposition~\ref{prop:registry} to a returned descriptor produces
its periodic class identifier. At future validation times apply the same
local construction to the new trajectory, compare its descriptor with
the stored prototype, and orient the endpoints using the recovered sign.
Choose the histogram rectangle strictly inside the protected rectangle.
For this rectangle, success is exactly the selected local branch, up to
the measured endpoint-coordinate errors: there are no missed protected
returns and no admitted foreign branch. Other trajectories remain failure
or outside-rectangle outcomes. There is no callback which recognizes
unknown lattice translates. A single-copy physical detector would not
implement this rule.

Only a rough gap estimate is needed to choose active validation times.
For example, when $|\widehat g-g|<d_*/100$, use
$T_1=2\widehat g+d_*/3$ and $T_2=2\widehat g+d_*/2$ on a sufficiently
small common rectangle. Both times lie strictly in the same collar.
The statistical inverse still recovers $g=W(0,0)/2$ and $A$ from the
observed densities; the rough probe value is not inserted as exact geometry.

\subsection{Square preparation and the discovery hazard}
Choose positions uniformly in $\Omega_L=[-L,L]^2$, reject solid positions,
and choose directions uniformly after a free position has been obtained.
Rejections are counted as launch attempts. Let $\tau_L$ be total variation
distance of the resulting position modulo the true lattice from quotient
free-position measure. For $L\ge2b$,
\begin{equation}\label{eq:square}
 \tau_L\le\min\{1,8V_1b/(A_0L)\},\qquad
 \Prob\{\hbox{a launch position is free}\}\ge\pi_0:=A_0/(4V_1).
\end{equation}
Indeed complete fundamental tiles inside the square cover its inner
square of half-width $L-b$. Their union $U$ has area at least
$4(L-b)^2$, and the omitted boundary layer has area at most $8Lb$.
Free positions in $U$ project exactly to the quotient law. The remaining
free positions are a contamination of weight at most
$8Lb/\{(A_0/V_1)4(L-b)^2\}$. This proves both inequalities without
requiring the apparatus to know any tile. A true periodic gate contracts
total variation; the finite-coordinate error will be treated separately.

Fix a finite time grid $\mathcal T$ in $[0,2R_-+d_*]$, fine enough
that for every $g\le R_-$ one grid point has
$T-2g\in[d_*/3,d_*/2]$. Write $H_0=|\mathcal T|$. Reduce the protected
rectangle to $[-a,a]^2$ so that $W-2g\le d_*/12$. The intrinsic
phase-density identity proved in Section~\ref{sec:inverse} gives, at
this grid point,
\[
 f_T(s,t)=\frac{-W_{st}}{2\pi A}(T-W)
                   \ge\frac{\omega_0d_*}{8\pi A_1}.
\]
Thus the quotient probability of a protected return is at least
\begin{equation}\label{eq:mu}
 \mu_*:=\min\left\{\tfrac12,
                         \frac{\omega_0d_*a^2}{2\pi A_1}\right\}>0.
\end{equation}
Choose a fixed discovery square large enough that $\tau_L\le\mu_*/2$.
At each proposal draw a fresh accepted free launch and choose its time
uniformly from $\mathcal T$. The local test and registry then discover
each still missing uniformly clear bridge class with conditional
probability at least
\begin{equation}\label{eq:hazard}
                         p_*=\frac{\mu_*}{2H_0}.
\end{equation}
This bound uses periodicity of the \emph{ideal geometric event}, whose
local implementation has just been proved; it does not assume a periodic
apparatus. Repeated records do not decrease the hazard.

\subsection{Why a finite witness exists}
\begin{lemma}[Saturated witness from short clear bridges]
\label{lem:witness}
The lifted graph of clear shortest bridges of length less than $R_-$ is
connected. Its finite quotient covers every body type and generates
$\Lambda$ by cycle translations. It contains a saturated witness of at
most $w_0$ records as in \eqref{eq:witness-main}.
\end{lemma}
\begin{proof}
The obstacle union has covering radius at most $b$: a translate of any
one obstacle point lies within a fundamental parallelogram diameter of
every prescribed point. Choose $b<c<R_-/2$. The open $c$-neighborhoods
of all bodies cover the connected plane, so their intersection graph is
connected. An intersecting pair has gap less than $R_-$. If its shortest
segment is blocked, any intervening body has a strictly smaller gap to
each endpoint body. Only finitely many pair-gap values below $R_-$ occur
modulo periods. Induction over these values replaces each blocked edge
by a path of shorter clear ones. This proves lifted connectivity.
A path from one body copy to any of its translates shows that quotient
closed walks generate every lattice vector.

Keep a quotient spanning tree, using at most $r_0-1$ edges. Two
independent fundamental cycles have determinant at most $B_0^2$, so
their subgroup has index at most $Q$ in $\Lambda$. Whenever the kept
subgroup is proper, some remaining fundamental cycle lies outside it.
Adjoining that one non-tree record makes the index a proper divisor,
hence at most half of the previous index. At most
$\lfloor\log_2 Q\rfloor$ further records suffice. Tree plus the two
initial non-tree edges gives \eqref{eq:witness-main}. No recorded pair
need be primitive. These bodies and paths occur only in the existence
proof, not in the acquisition algorithm's input.
\end{proof}

If $K_{\rm cov}$ is the first proposal epoch containing the witness,
iterated conditioning and a union bound give
\begin{equation}\label{eq:coupon}
 \Prob\{K_{\rm cov}>k\}\le w_0(1-p_*)^k\le w_0e^{-p_*k}.
\end{equation}
Independence between the discoveries of different records is unnecessary.
The conditional bound comes from fresh launches, not from an assumed
hazard of a complete-geometry oracle. The same local registry also accepts
any other genuine producer; its more general hazard theorem is retained
in the preceding manuscript.
